\begin{helper}[No middle-trunk escape at height five]
Let $N$ be normal globally maximum data for a $k$-triple
$(T,e_0,\varphi)$, suppose that $N.\mathsf{height}=h(G)=5$, and write
$y_i=T.\mathsf{vertex}(i)$ for $0\leq i\leq4$.  Let
$\gamma\in\overline\varphi(y_0)$.  Suppose that an exact edge identity
$g$ has color $\gamma$ and endpoints $\{y_i,v\}$ for some
$i\in\{0,1,2,3,4\}$, where
$v\in V(T)\setminus\{y_0,y_1,y_2,y_3,y_4\}$.
If another exact edge identity $r$ has color $\gamma$ and endpoints
$\{y_2,w\}$ or $\{y_3,w\}$, then
\[
 w\in\{y_0,y_1,y_2,y_3,y_4\}.
\]
All assertions concern the named multigraph edge identities, so no two
parallel edges are identified merely because their unordered endpoint
pairs agree.
\end{helper}
\begin{proof}
Put
\[
 a_1=T.\mathsf{edge}(1),\qquad
 a_2=T.\mathsf{edge}(2),\qquad
 a_3=T.\mathsf{edge}(3),\qquad q=N.\mathsf{closingEdge}.
\]
Thus $a_1,a_2,a_3$ are the paper's identities $e_2,e_3,e_4$.
The height upper bound, closing-index equation, endpoint fields, and
support fields of \noderef{NormalGlobalMaximumTriple}, together with
the first-edge, attachment, and injectivity fields of
\noderef{TreeSequence} and the endpoint meaning of
\noderef{Incident}, give
\[
 \operatorname{ends}(a_1)=\{y_1,y_2\},\quad
 \operatorname{ends}(a_2)=\{y_2,y_3\},\quad
 \operatorname{ends}(a_3)=\{y_3,y_4\},\quad
 \operatorname{ends}(q)=\{y_4,y_0\}.                 \tag{1}
\]
The five $y_i$ are pairwise distinct, and the stored closing path is the
simple incidence path
\[
 y_1,a_1,y_2,a_2,y_3,a_3,y_4,q,y_0.                 \tag{2}
\]

We next recover all four colors in (2), rather than treating alternation
as implicit.  Every edge-node on (2) has color $N.\mathsf{alpha}$ or
$N.\mathsf{beta}$ by \noderef{TwoColorIncidenceGraph}.  The color
$N.\mathsf{beta}$ is missing at $y_1$ by
\noderef{NormalGlobalMaximumTriple}; expanding
\noderef{MissingColors} and \noderef{PresentColors}, the incident identity
$a_1$ therefore cannot have that color.  Hence
$\varphi(a_1)=N.\mathsf{alpha}$.  The global field of $N$, expanded by
\noderef{GlobalMaximumTriple} and \noderef{KTriple}, makes $\varphi$
proper in the sense of \noderef{IsProperEdgeColoring}.  Consecutive
identities in (2) are distinct and adjacent in the sense of
\noderef{AdjacentEdges}, so properness alternates their colors.
Consequently
\[
 \varphi(a_1)=\varphi(a_3)=N.\mathsf{alpha},\qquad
 \varphi(a_2)=\varphi(q)=N.\mathsf{beta}.            \tag{3}
\]
In particular, the incidence of $a_1$ at $y_1$ also proves
$N.\mathsf{alpha}\ne N.\mathsf{beta}$.

Apply \noderef{ScheideKempePathHeightBound} to $N.\mathsf{global}$,
$\gamma$ missing at $y_0$, and $N.\mathsf{beta}$ missing at $y_1$.
It supplies a simple incidence path $P_1$ from $y_0$ to $y_1$ whose
support is the full $(\gamma,N.\mathsf{beta})$-component $C_1$ in the
sense of \noderef{TwoColorComponent}, with
\[
 4\mid\operatorname{length}(P_1),\qquad
 \operatorname{length}(P_1)/2+1\leq h(G)=5.          \tag{4}
\]
By (1)--(3), the incidences $y_0,q,y_4$ lie in $C_1$.  Also
\noderef{ScheideGlobalMaximumOrder}, read through
\noderef{ColoredVertexSetElementary}, makes $V(T)$ elementary.  Since
$\gamma$ is missing at $y_0$ and $y_0\ne y_1$, it is not missing at
$y_1$.  The definitions \noderef{MissingColors} and
\noderef{PresentColors} therefore give a $\gamma$-colored exact identity
$b$ incident with $y_1$; its edge-node lies in $C_1$.

We record the two required color separations.  First
$\gamma\ne N.\mathsf{beta}$, since $b$ has color $\gamma$ at $y_1$ while
$N.\mathsf{beta}$ is missing there.  Suppose instead that
$\gamma=N.\mathsf{alpha}$.  Then (1)--(3) put all five trunk vertices,
and the four exact edge-nodes $q,a_1,a_2,a_3$, in $C_1$.  The specified
identity $g$ joins one of those vertices to $v$, so $v$ also lies in
$C_1$.  These six vertex-nodes are distinct by the hypotheses on $v$;
the four edge-nodes are distinct by edge injectivity and their exact
endpoint pairs in (1).  Thus $P_1$ has at least ten distinct incidence
nodes.  Simplicity of $P_1$ gives
$\operatorname{length}(P_1)\geq9$; divisibility by four then gives length
at least twelve, contradicting (4).  Hence
\[
 \gamma\ne N.\mathsf{alpha},\qquad
 \gamma\ne N.\mathsf{beta}.                         \tag{5}
\]
This count retains the edge-nodes $q,a_1,a_2,a_3$ as identities and makes
no simplicity assumption about the underlying multigraph.

Assume for a contradiction that $w$ is distinct from every trunk vertex.
We claim that none of $y_2,y_3,w$ belongs to $C_1$.  If any one belonged,
the $N.\mathsf{beta}$-identity $a_2$ between $y_2,y_3$ and the
$\gamma$-identity $r$ between $w$ and one of those two vertices would put
all three in $C_1$.  Then $C_1$ would contain the six distinct vertex-nodes
$y_0,y_1,y_2,y_3,y_4,w$ and the four edge-nodes $q,a_2,b,r$.
These edge identities are distinct: $q\ne a_2$ by their endpoint pairs in
(1), the cross-color pairs are distinct by (5), and $b\ne r$ because $b$
is incident with $y_1$ whereas neither possible endpoint pair of $r$
contains $y_1$.  As before, ten distinct incidence nodes force path length
at least nine and then at least twelve, contrary to (4).  Therefore
\[
 y_2,y_3,w\notin C_1,\qquad \operatorname{inr}(r)\notin C_1.             \tag{6}
\]
For the final clause, if the edge-node of $r$ belonged to $C_1$, its
incidence adjacency in \noderef{TwoColorIncidenceGraph} would put its
middle-trunk endpoint in $C_1$, contrary to the first part of (6).

We now use the permitted switch, checking its hypotheses explicitly.
By \noderef{ScheideGlobalMaximumFixedPairMaximal}, the globally maximum
triple has $T$ fixed-pair maximal.  Apply
\noderef{ScheideMaximalKempePath} at tree indices zero and one with colors
$\gamma,N.\mathsf{beta}$.  The premise
\noderef{TreeColorUnusedThrough} is vacuous at index one, since it asks
only about colored tree-edge indices $j$ satisfying $1\leq j<1$.
Consequently the component interchange
\[
 \varphi'=\noderef{TwoColorComponentSwap}
   (G,\varphi,\gamma,N.\mathsf{beta},y_0)
\]
is proper, still leaves exactly $e_0$ uncolored, and leaves the same
ordered sequence $T$ Tashkinov.  These are precisely the color-dependent
clauses of \noderef{KTriple}, \noderef{IsOneEdgeUncolored}, and
\noderef{TashkinovTree}; the graph-theoretic clauses and the cardinality
of $V(T)$ do not change.  Hence $(T,e_0,\varphi')$ again satisfies
\noderef{GlobalMaximumTriple}.

The missing-color conclusions after the switch follow clause by clause
from \noderef{TwoColorComponentSwapMissingBehavior}.  The endpoint $y_1$
belongs to $C_1$ and formerly missed $N.\mathsf{beta}$, so it misses
$\gamma$ under $\varphi'$.  By (5), the unaffected color
$N.\mathsf{alpha}$ remains missing at $y_0$.  The definition of the swap,
together with (3) and (6), gives
\[
 \varphi'(q)=\gamma,\quad
 \varphi'(a_1)=\varphi'(a_3)=N.\mathsf{alpha},\quad
 \varphi'(r)=\gamma.                                \tag{7}
\]
Here $q$ belongs to $C_1$, the two alpha identities are fixed because of
(5), and $r$ is fixed because its exact edge-node is outside $C_1$.

Apply \noderef{ScheideKempePathHeightBound} a second time to the switched
globally maximum triple, using $N.\mathsf{alpha}$ missing at $y_0$ and
$\gamma$ missing at $y_1$.  Its full
$(N.\mathsf{alpha},\gamma)$-component has a simple incidence path $P_2$
with length divisible by four and
\[
 \operatorname{length}(P_2)/2+1\leq h(G)=5.          \tag{8}
\]
Equations (1) and (7) force the segments
\[
 y_0,q,y_4,a_3,y_3,qquad y_1,a_1,y_2
\]
into this component.  The exact identity $r$ then brings in $w$, whether
its middle endpoint is $y_2$ or $y_3$.  Thus $P_2$ contains the six
distinct vertex-nodes $y_0,y_1,y_2,y_3,y_4,w$ and the four exact
edge-nodes $q,a_1,a_3,r$.  The four edge identities are distinct: (1) and
tree edge injectivity separate $q,a_1,a_3$, (5) separates $r$ from the
two alpha identities, and the old colors in (3) separate $r$ from $q$.
The same ten-node support count forces length at least twelve, contradicting
(8).  This contradiction discharges the assumption that $w$ lies outside
the trunk, and proves the asserted five-way alternative.
\end{proof}
